\begin{appendices}

\section{Demonstração manual do dimensionamento}\label{ap:demonstracao}

\begin{fillblock}
\noindent\textbf{[A PREENCHER]} Apresentar, passo a passo e com base no equacionamento da Seção~\ref{sec:dimensionamento}, o dimensionamento manual completo de uma longarina circular e de uma prancha do tabuleiro, para uma das soluções da fronteira eficiente selecionada nos casos de estudo da Seção~\ref{sec:casos} (recomenda-se uma solução intermediária da Ponte~01, com valores de $d$, $b_w$, $h$, $esp_{long}$ e $esp_{tab}$ definidos após a otimização final). Sequência sugerida, compatível com a rotina \textsc{VerificaEstrutura} do Algoritmo~\ref{alg:nsga2}:

\begin{enumerate}
    \item Dados de entrada da solução selecionada: $d$, $b_w$, $h$, $esp_{long}$, $esp_{tab}$, vão $L$, trem tipo, propriedades da madeira (Tabelas~\ref{tab:parametros} e~\ref{tab:casos});
    \item Propriedades geométricas das seções da longarina e do tabuleiro (Equações~\eqref{eq:areacirc} a~\eqref{eq:iret});
    \item Ações permanentes (peso próprio da longarina e do tabuleiro) e ações variáveis do trem tipo, incluindo o coeficiente de impacto vertical (Equações~\eqref{eq:permanente} a~\eqref{eq:civ2});
    \item Esforços solicitantes de momento fletor, esforço cortante e flecha, para as ações permanente e variável (Equações~\eqref{eq:mgk} a~\eqref{eq:flechaq});
    \item Combinação última de ações e resistências de cálculo, com o coeficiente de modificação $k_{mod}$ (Equações~\eqref{eq:comb} e~\eqref{eq:xd});
    \item Verificação da longarina à flexão (Equação~\eqref{eq:verflexao}), ao cisalhamento (Equação~\eqref{eq:vercis}) e à flecha, e verificação do tabuleiro à flexão;
    \item Síntese das funções limite $g_1$ a $g_4$ obtidas manualmente, comparando-as com a saída da plataforma computacional para a mesma solução.
\end{enumerate}
\end{fillblock}

\section{Relatório de verificação dos elementos}\label{ap:relatorio}

A plataforma computacional emite, ao final da etapa de projeto detalhado, um relatório completo contendo as propriedades geométricas das seções, os esforços solicitantes (momentos fletores, esforços cortantes e flechas), as combinações de ações e o resultado de cada verificação de estado limite último e de serviço para a longarina e para o tabuleiro. O relatório reproduz, para a solução selecionada pelo projetista, a sequência de cálculo demonstrada no Apêndice~\ref{ap:demonstracao}.

\begin{fillblock}
\noindent\textbf{[A PREENCHER]} Anexar o relatório completo (exportado pela plataforma) para a mesma solução do Apêndice~\ref{ap:demonstracao}, incluindo a tabela de utilização (grau de solicitação) de cada verificação -- ver a Figura~\ref{fig:utilizacao} da Seção~\ref{sec:ponte01}.
\end{fillblock}

\end{appendices}
